% Part: proof-theory
% Chapter: natural-deduction
% Section: grafting

\documentclass[../../../include/open-logic-section]{subfiles}

\begin{document}

\olfileid{pt}{nat}{gra}

\olsection{Nyambung \usetoken{P}{proof}}

!!^{proof}s ing \Log{N1c} lan~\Log{N1i} iku !!{proof}s
kanggo !!{formula}s~$!A$ saka asumsi~$!B$.
Anggepen $!B$ dhéwé nduwèni !!a{proof}.
Kepriyé nggabungaké !!{proof} kanggo~$!B$ karo
!!{proof} kanggo $!A$ saka~$!B$, supaya ngolèhké
!!a{proof} kanggo~$!A$ tanpa nganggo asumsi~$!B$?

Salah siji cara yaiku mbusak asumsi~$!B$ saka
!!{proof} kanggo~$!A$ kanthi ngetrapaké~\Intro{\lif}.
Banjur !!{proof} kang diasilaké kanggo~$!B \lif !A$
bisa digandhèng karo !!{proof} kanggo~$!B$ saéngga dadi
\[
\AxiomC{$\Discharge{!B}{x}$}
\DeduceC{$!A$}
\DischargeRule{\Intro{\lif}}{x}
\UnaryInfC{$!B \lif !A$}
\AxiomC{}
\DeduceC{$!B$}
\RightLabel{\Elim{\lif}}
\BinaryInfC{$!A$}
\DisplayProof
\]
Cara iki gampang, nanging rada ora éndah: kena apa
nglebokaké $\lif$ mung kanggo langsung ngilangaké?
Liku-liku liwat $!B \lif !A$ iki kuduné bisa diendhani.
(Pancen, teorema normalisasi mratélakaké yèn
liku-liku mangkono mesthi bisa diendhani.)

Ing \Log{N1c} lan~\Log{N1i}, liku-liku iki gampang
diendhani: saben asumsi~$!B^x$ cukup diganti
sakabèhé !!{proof} kanggo~$!B$.
``Nyambung'' !!{proof}s ing papan asumsi bukti liya
iki migunani, mula kita tetepaké dadi definisi.

\begin{defn}
Yèn $\delta_1$ iku \Log{N1c}-!!{proof} (utawa
\Log{N1i}-!!{proof}) kanggo~$!A$ saka asumsi
kabuka~$!B^x$, lan $\delta_2$ iku \Log{N1c}-!!{proof}
(\Log{N1i}-!!{proof}) kanggo~$!B$, banjur
$\Subst{\delta_1}{\delta_2}{!B^x}$ yaiku asil saka
ngganti saben kedadéan~$!B^x$ kang durung dibusak
ing~$\delta_1$ nganggo~$\delta_2$.
\end{defn}

Yèn $\delta_1$ lan $\delta_2$ ora nduwèni
!!{formula}s asumsi béda kanthi label kang padha,
lan ora ana asumsi kabuka ing~$\delta_2$ kang ngemot
variabel eigen saka~$\delta_1$, banjur
$\Subst{\delta_1}{\delta_2}{!B^x}$ iku !!{proof}
kang bener, kanthi asumsi kabuka saka $\delta_1$
lan saka~$\delta_2$. Nalika nyambung !!{proof}s,
syarat iki lumrahé kacukupan. Syarat kapisan kacukupan
yèn $\delta_1$ lan $\delta_2$ subbukti saka !!{proof}
kang luwih gedhé. Yèn syarat kapindho durung kacukupan,
variabel eigen ing~$\delta_1$ bisa diganti jenengé
supaya syarat mau kacukupan.

Definisi kang cocog uga lumaku kanggo \Log{N2c} lan~\Log{N2i}.

\begin{defn}
Yèn $\delta_1$ iku \Log{N2c}-!!{proof} (utawa
\Log{N2i}-!!{proof}) kanggo~$x:!B, \Gamma_1 \Sequent !A$,
lan $\delta_2$ iku \Log{N2c}-!!{proof}
(\Log{N2i}-!!{proof}) kanggo~$\Gamma_2 \Sequent !B$,
banjur $\Subst{\delta_1}{\delta_2}{x:!B}$ iku asil
saka ngganti saben aksioma $x:!B \Sequent !B$
ing~$\delta_1$ nganggo~$\delta_2$, lan nambah
$\Gamma_2$ ing kontèks saben sekèn ing sangisoré
aksioma mau.
\end{defn}

\begin{prop}
Yèn $\delta_1 \Proves x:!B, \Gamma_1 \Sequent !A$
lan $\delta_2 \Proves \Gamma_2 \Sequent !B$,
banjur $\Subst{\delta_1}{\delta_2}{x:!B}
\Proves \Gamma_1, \Gamma_2 \Sequent !A$,
anggèr ora ana variabel eigen saka~$\delta_1$
ing~$\Gamma_2$.
\end{prop}


\begin{proof}
  Kanthi induksi tumrap $\pheight{\delta_1}$.
\end{proof}

\end{document}
